% ============================================================
\chapter{Stationarity and Autocovariance}
\label{ch:stationnarite}
% ============================================================

\section{Motivating example}
The daily temperature in Paris over 30 years exhibits a pronounced seasonality, but after deseasonalisation, the residuals appear to fluctuate around a constant mean with a time-invariant dependence structure. This is \textbf{stationarity}.

\section{Strict and weak stationarity}

\begin{definition}[Strict stationarity]\index{stationarity!strict}
$(X_t)$ is \textbf{strictly stationary} if for all $k\geq 1$, all $(t_1,\ldots,t_k)$, and all $h\in\Z$:
\[
(X_{t_1},\ldots,X_{t_k}) \stackrel{d}{=} (X_{t_1+h},\ldots,X_{t_k+h}).
\]
\end{definition}

\begin{definition}[Weak stationarity (second-order)]\index{stationarity!weak}
$(X_t)$ is \textbf{weakly stationary} (or wide-sense stationary) if:
\begin{enumerate}
\item $\E[X_t^2]<\infty$ for all $t$,
\item $\mu = \E[X_t]$ does not depend on $t$,
\item $\gamma(t,t+h) = \Cov(X_t,X_{t+h})$ depends only on $h$.
\end{enumerate}
\end{definition}

\begin{remark}
Strict stationarity + finite second moments $\Rightarrow$ weak stationarity. The converse is false in general, but holds for Gaussian processes.
\end{remark}

\section{Autocovariance function}

\begin{definition}[Autocovariance and autocorrelation]\index{autocovariance}\index{autocorrelation}
For a weakly stationary process:
\begin{align}
\gamma(h) &= \Cov(X_t,X_{t+h}) = \E[(X_t-\mu)(X_{t+h}-\mu)], \\
\rho(h) &= \frac{\gamma(h)}{\gamma(0)} = \Corr(X_t,X_{t+h}).
\end{align}
The function $\rho(\cdot)$ is the \textbf{autocorrelation function} (ACF)\index{ACF}.
\end{definition}

\begin{theorem}[Properties of $\gamma$]\label{thm:gamma-prop}
\begin{enumerate}
\item $\gamma(0)\geq 0$ (it is the variance).
\item $\gamma(h)=\gamma(-h)$ (symmetry).
\item $|\gamma(h)|\leq\gamma(0)$ (by Cauchy--Schwarz).
\item $\gamma$ is \textbf{positive semi-definite}: for all $n$, all $(a_1,\ldots,a_n)\in\R^n$, and all $(t_1,\ldots,t_n)$:
\[
\sum_{i,j}a_i a_j\gamma(t_i-t_j)\geq 0.
\]
\end{enumerate}
\end{theorem}

\begin{proof}
(4): $\sum_{i,j}a_ia_j\gamma(t_i-t_j)=\Var\!\left(\sum_i a_i X_{t_i}\right)\geq 0$.
\end{proof}

\begin{theorem}[Herglotz]\index{Herglotz}
$\gamma:\Z\to\R$ is an autocovariance function if and only if it is positive semi-definite. Moreover, there exists a finite positive measure $F$ on $[-\pi,\pi]$ (spectral measure) such that:
\[
\gamma(h) = \int_{-\pi}^{\pi} e^{i\omega h}\,dF(\omega).
\]
\end{theorem}

\section{Estimation of the autocovariance}

\begin{definition}[Sample autocovariance]\index{autocovariance!sample}
\[
\hat{\gamma}(h) = \frac{1}{n}\sum_{t=1}^{n-|h|}(X_t-\bar{X})(X_{t+|h|}-\bar{X}), \quad \hat{\rho}(h)=\frac{\hat{\gamma}(h)}{\hat{\gamma}(0)}.
\]
Dividing by $n$ (rather than $n-|h|$) ensures that the matrix $(\hat{\gamma}(i-j))$ is positive semi-definite.
\end{definition}

\begin{theorem}[Bartlett]\index{Bartlett}
For a white noise of size $n$: $\hat{\rho}(h)\stackrel{\text{approx}}{\sim}\mathcal{N}(0,1/n)$ for $h\geq 1$. The \textbf{Bartlett bands} $\pm 1.96/\sqrt{n}$ are used to test whether $\rho(h)=0$.
\end{theorem}

\section{Partial autocorrelation}\index{PACF}

\begin{definition}[PACF]
The \textbf{partial autocorrelation} at lag $h$ is:
\[
\alpha(h) = \Corr(X_t-\hat{X}_t,\;X_{t+h}-\hat{X}_{t+h}),
\]
where $\hat{X}_t$ and $\hat{X}_{t+h}$ are the linear projections onto $(X_{t+1},\ldots,X_{t+h-1})$.

Equivalently, $\alpha(h)$ is the last coefficient $\phi_{hh}$ in the linear regression of $X_t$ on $(X_{t-1},\ldots,X_{t-h})$.
\end{definition}

\begin{theorem}
For an AR$(p)$: $\alpha(h)=0$ for $h>p$. This is the main tool for identifying the order $p$.
\end{theorem}

\section{Stationarity tests}\index{test!Dickey--Fuller}\index{test!KPSS}

\begin{definition}[Augmented Dickey--Fuller test (ADF)]
To test for the presence of a unit root in $X_t = \phi X_{t-1}+\eps_t$:
\[
H_0 : \phi=1 \quad\text{(non-stationary)} \quad \text{vs} \quad H_1 : |\phi|<1 \quad\text{(stationary)}.
\]
One estimates $\nabla X_t = (\phi-1)X_{t-1}+\eps_t$ and tests whether $\phi-1<0$.
\end{definition}

\begin{warning}
The distribution of the ADF statistic is \textbf{not} Gaussian under $H_0$. One must use the Dickey--Fuller tables (or critical values computed by simulation).
\end{warning}

\section{Implementation}

\begin{codeblock}[title=\textbf{Python}]
\begin{minted}{python}
import numpy as np
import matplotlib.pyplot as plt
from statsmodels.tsa.stattools import acf, pacf, adfuller
from statsmodels.graphics.tsaplots import plot_acf, plot_pacf

# Simuler un AR(1)
np.random.seed(0)
n = 500
phi = 0.7
eps = np.random.normal(0, 1, n)
x = np.zeros(n)
for t in range(1, n):
    x[t] = phi * x[t-1] + eps[t]

# ACF et PACF
fig, axes = plt.subplots(1, 2, figsize=(12, 4))
plot_acf(x, lags=30, ax=axes[0], title='ACF')
plot_pacf(x, lags=30, ax=axes[1], title='PACF')
plt.tight_layout()
plt.savefig('ch02_acf_pacf.pdf')
plt.show()

# Test ADF
result = adfuller(x, regression='c')
print(f"ADF statistic: {result[0]:.4f}")
print(f"p-value: {result[1]:.4f}")
print(f"Stationnaire: {result[1] < 0.05}")

# Marche aléatoire (non stationnaire)
rw = np.cumsum(eps)
result_rw = adfuller(rw, regression='c')
print(f"\nMarche aléatoire - ADF: {result_rw[0]:.4f}, p={result_rw[1]:.4f}")
\end{minted}
\end{codeblock}

\section{Exercises}

\begin{exercise}[$\star$]
Compute the autocovariance function of a white noise $\eps_t\sim\mathrm{WN}(0,\sigma^2)$.
\end{exercise}

\begin{exercise}[$\star$]
Compute the ACF of $X_t = \eps_t + \theta\eps_{t-1}$ (MA(1)). Show that $\rho(h)=0$ for $|h|\geq 2$.
\end{exercise}

\begin{exercise}[$\star\star$]
Show that if $(X_t)$ is strictly stationary with $\E[X_t^2]<\infty$, then $(X_t)$ is weakly stationary.
\end{exercise}

\begin{exercise}[$\star\star$]
Prove that the sample autocovariance $\hat{\gamma}(h)$ (divided by $n$) yields a positive semi-definite Toeplitz matrix.
\end{exercise}

\begin{exercise}[$\star\star\star$ -- Project]
Apply the ADF test and the KPSS test to 10 economic time series (GDP, unemployment, inflation, etc.). Compare the conclusions and discuss cases of disagreement.
\end{exercise}

\begin{keyformulas}
\begin{align*}
\gamma(h) &= \Cov(X_t,X_{t+h}), \quad \rho(h)=\gamma(h)/\gamma(0) \\
\hat{\rho}(h) &\stackrel{\text{approx}}{\sim} \mathcal{N}(0,1/n) \quad\text{under }H_0:\rho(h)=0 \\
\text{ADF} &: \nabla X_t = (\phi-1)X_{t-1}+\eps_t, \quad H_0:\phi=1
\end{align*}
\end{keyformulas}
